\documentclass[fontsize=12 pt]{scrartcl}
\usepackage{setspace}
\onehalfspacing
\usepackage{amsmath,amssymb,amsfonts,amsthm,mathtools}
\usepackage{algorithm}
\usepackage{algpseudocode}
\usepackage[english]{babel}
\usepackage[T1]{fontenc}
\usepackage[utf8x]{inputenc}
\usepackage{lmodern}
\usepackage{dsfont}
\usepackage{float}
\usepackage{bbm}
\usepackage[round]{natbib}
\usepackage{color} 
\usepackage[defaultlines=2,all]{nowidow}
\usepackage{caption}
\usepackage[labelformat=simple]{subcaption}
\renewcommand\thesubfigure{(\alph{subfigure})}

\setlength\parindent{0pt}
\setlength{\parskip}{6pt plus 1pt minus 1pt}

\newcommand{\red}{\textcolor{red}}


\begin{document}

% --- TITLE PAGE (Unchanged) ---
\begin{titlepage}
    \centering
    {\scshape\LARGE TU Dortmund \par}
    \vspace{1cm}
    {\scshape\Large Fachgebiet Statistical Methods for Big Data \par}
    \vspace{1cm}
    {\huge\bfseries Advanced Component-wise Gradient Boosting for Improved Generalisability\par}
        {\Large \par}
    \vspace{2cm}
    {\Large Author:\\
            André Kafanke \par}
    \vspace{3cm}
        {\Large Reviewer:\\
        Prof. Dr. Andreas Groll\\\par}    
        {\Large Prof. Dr. Andreas Mayr\\\par}    
    \vfill
    {\large March 02, 2026\par}
\end{titlepage}

\tableofcontents
\newpage
\listoffigures
\cleardoublepage

% --- CONTENT START ---

\section{Introduction}
% Motivation: Overfitting in Boosting, Generalization gap, Need for better feature selection in high-dim.

\newpage
\section{Statistical Methods}

\subsection{Supervised Learning Foundations}
\subsubsection{The Regression Setting and the I.I.D. Assumption}

In supervised learning, specifically in the regression setting, we consider an approximation problem of a mapping function $f: \mathcal{X} \rightarrow \mathcal{Y}$ from an input space $\mathcal{X} \subseteq \mathbb{R}^p$ to an output space $\mathcal{Y} \subseteq \mathbb{R}$, where $p$ denotes the dimensionality of the feature vector. \citep{hastie2009elements}. The main objective is to learn the relationship between the input and target variables resulting in an approximation of the mapping function. 

The data generation process is assumed to be characterized by an unknown joint probability distribution $P(\mathbf{x}, y)$, which is defined in the product space $\mathcal{X} \times \mathcal{Y}$. This probability distribution captures the conditional distribution of the target $P(y|\mathbf{x})$ as well as the marginal distribution of the input features $P(\mathbf{x})$\citep{bishop2006pattern}. We can formalize this relationship as:

\begin{equation}
    y = f^*(\mathbf{x}) + \epsilon
\end{equation}

where $f^*(\mathbf{x}) = \mathbb{E}[y|\mathbf{x}]$ describes the true target function and $\epsilon$ represents a random noise term, assumed to be independent of $\mathbf{x}$ and normally distributed with mean zero and variance $\sigma^2$, denoted by $\epsilon \sim \mathcal{N}(0, \sigma^2)$ \citep{hastie2009elements}.

We are provided with a dataset $\mathcal{D} = \{(\mathbf{x}_i, y_i)\}_{i=1}^n$, containing $n$ observations. One of the fundamental assumptions in statistical learning theory is that these observations are drawn \textit{independently and identically distributed} (i.i.d.) from the joint distribution $P(\mathbf{x}, y)$ \citep{vapnik1998statistical}.

The \textbf{independence} between observations implies that one realization of an observation $(\mathbf{x}_i, y_i)$ yields no information about another realization of an observation $(\mathbf{x}_j, y_j)$ for $i \neq j$. The assumption of \textbf{identical distribution} implies the data generation process $P(\mathbf{x}, y)$ remains stationary during data collection and between training and testing \citep{mohri2018foundations}. 

Under i.i.d. assumptions, we can approximate the expected risk $R(f)$ with the empirical risk $R_{emp}(f)$ calculated over the training data:

\begin{equation}
    R(f) = \int_{\mathcal{X} \times \mathcal{Y}} L(y, f(\mathbf{x})) \, dP(\mathbf{x}, y) \approx \frac{1}{n} \sum_{i=1}^n L(y_i, f(\mathbf{x}_i))
\end{equation}

where $L(\cdot, \cdot)$ denotes a loss function, such as the squared error $L(y, \hat{y}) = (y - \hat{y})^2$, which is used in this work primarily. In theory, convergence of the empirical to the expected risk with $n \to \infty$ is derived from the i.i.d assumptions, but in real-world scenarios they are often violated through circumstances such as concept drifts (changes of $P(y|\mathbf{x})$) or covariate shifts (changes of $P(\mathbf{x})$) \citep{quionero2009dataset}. This highlights the need for robust regularization and feature selection methods, which are introduced and investigated later in this thesis. 

\subsubsection{Loss Functions and Empirical Risk Minimization}

To evaluate the quality of a function $f \in \mathcal{F}$, we need a non-negative loss function $L: \mathcal{Y} \times \mathcal{Y} \rightarrow \mathbb{R}_{\geq 0}$ which quantifies the divergence from the predicted outcome $\hat{y} = f(\mathbf{x})$ to the true outcome $y$ \citep{hastie2009elements}.

Since in this thesis, we focus on regression problems with additive Gaussian noise, we primarily utilize the Squared Error loss ($L_2$ loss). We define the loss for a single observation as:

\begin{equation}
    L(y, f(\mathbf{x})) = \frac{1}{2} (y - f(\mathbf{x}))^2
\end{equation}

 Assuming the target variable follows a normal distribution conditional on the input variables ($y | \mathbf{x} \sim \mathcal{N}(f(\mathbf{x}), \sigma^2)$), then minimizing this loss function over the training dataset gives the Maximum Lielihood Estimate (MLE) for the model parameters \citep{bishop2006pattern}.

The expected risk $R(f)$ cannot be minimized analytically, since we have no information about the true joint distribution $P(\mathbf{x}, y)$. Therefore, we utilize the Empirical Risk Minimization (ERM) principle presented by \citet{vapnik1998statistical}. It states that a learning algorithm should select the function $\hat{f}$ which minimizes the average loss over a training sample $\mathcal{D}$:

\begin{equation}
    \hat{f} = \underset{f \in \mathcal{F}}{\arg \min} \mathcal{R}_{emp}(f) = \underset{f \in \mathcal{F}}{\arg \min} \frac{1}{n} \sum_{i=1}^n L(y_i, f(\mathbf{x}_i))
\end{equation}

While the ERM principle is a powerful induction tool, it needs good conditions to yield qualitative results. In cases where the hypothesis space $\mathcal{F}$ is very complex, the dimensionality $p$ is high relative to the sample size $n$, or the noise $\epsilon$ is very high relative to the target function $f^*(\mathbf{x})$, the ERM becomes prone to \textit{overfitting}. This involves memorization of the noise $\epsilon$ instead of learning the target function $f^*(\mathbf{x})$, ultimately leading to poor predictive quality on unseen data, i.e. poor generalization \citep{mohri2018foundations}.

To mitigate the risk of overfitting, modern machine learning algorithms offer a variety of methods, but generally either extend the learnable function with a regularization term or induce implicit regularization mechanisms. In later sections we will show that in Component-wise Gradient Boosting, implicit regularization is given via the number of boosting iterations and shrinkage of updates \citep{zhang2005boosting}. In contrast, explicit regularization methods such as \textit{Flooding} (Section 2.5.2) directly prevent the empirical risk from reaching zero \citep{ishida2021flooding}.

\subsubsection{Cross-Validation Strategies}
\label{sec:cv}

To evaluate the predictive performance of the Componentwise Gradient Boosting model and to determine an optimal model, we employ a nested validation strategy. It utilizes distinct methods for \textit{model selection} and for \textit{model evaluation}, which have the purpose of preventing optimism bias and data leakage \citep{cawley2010over}.

\paragraph{Inner Loop: Model Selection via K-Fold Cross-Validation}

The complexity of a boosting model is primarily regularized over the number of iterations of fitting $m_{stop}$. While a too small number of iterations is associated with insufficient capacity to capture the true target function (underfitting), a too large number can lead to the model fitting the noise (overfitting) \citep{buhlmann2007boosting}. To estimate an optimal stopping point $m_{stop}$, we apply K-Fold Cross-Validation on the training data $\mathcal{D}_{train}$.

Let $\mathcal{S}_1, \dots, \mathcal{S}_K$ denote $K$ disjoint subsets of equal size, such that $\mathcal{D}_{train} = \bigcup_{k=1}^K \mathcal{S}_k$ \citep{kohavi1995study}. The algorithm then iterates through $k = 1, \dots, K$, where for each iteration:

\begin{enumerate}
    \item A subsidiary model $\hat{f}^{k}$ is fitted on the set $\mathcal{D}_{train} \setminus \mathcal{S}_k$.
    \item The empirical risk (validation loss) is calculated on the held-out validation fold $\mathcal{D}_{val} = \mathcal{S}_k$ for each iteration $m = 1, \dots, M_{max}$.
\end{enumerate}

Let $L(y, \hat{f}_m(\mathbf{x}))$ be the loss function at iteration $m$. The cross-validation loss is retained by the average of the empirical risk over all K folds.

\begin{equation}
    \text{CV}(m) = \frac{1}{K} \sum_{k=1}^K \left( \frac{1}{|\mathcal{S}_k|} \sum_{(\mathbf{x}_i, y_i) \in \mathcal{S}_k} L(y_i, \hat{f}^{k}_m(\mathbf{x}_i)) \right)
\end{equation}

We define the optimal iteration $m^*$ as the index, which minimizes the cross-validation loss \citep{hastie2009elements}:

\begin{equation}
    m^* = \underset{m \in \{1, \dots, M_{max}\}}{\arg \min} \text{CV}(m)
\end{equation}

After determining $m^*$, the final model $\hat{f}_{m^*}$ is obtained by fitting it on the full training set $\mathcal{D}_{train}$, while stopping at iteration $m^*$. This leads to the model utilizing all available training observations, while simultaneously being restricted in complexity by cross-validation.

\paragraph{Outer Loop: Model evaluation via Monte Carlo Cross-Validation}

To ensure an realistic evaluation of the generalisation capabilities of a chosen set of model hyperparameters, we apply Monte Carlo Cross-Validation, also known as Repeated Random Subsampling. Therefore, we nest the inner K-fold selection loop within the outer evaluation loop \citep{arlot2010survey}.

For each training run of the experiment, the full dataset $\mathcal{D}$ is randomly split into a training $\mathcal{D}_{train}$ and test set $\mathcal{D}_{test}$. Then, the model selection is performed on the training set $\mathcal{D}_{train}$, while the test set is held out of the process to then be used to compute the generalization error on the final model $\hat{f}_{m^*}$:

\begin{equation}
    \hat{R}_{gen} = \frac{1}{|\mathcal{D}_{test}|} \sum_{(\mathbf{x}_i, y_i) \in \mathcal{D}_{test}} L(y_i, \hat{f}_{m^*}(\mathbf{x}_i))
\end{equation}

This procedure is repeated over multiple independent trails with random seeds, which allows for estimation of both the expected prediction error and the variance of prediction errors stemming from sampling variability \citep{dietterich1998approximate}.

\subsubsection{Distributional Shifts}

As previously stated, the ERM principle is highly dependent on the assumption that the training data $\mathcal{D}_{train}$ and the test data $\mathcal{D}_{test}$ are drawn from the same joint distribution $P(\mathbf{x}, y)$. In many real-world scenarios, however, this assumption cannot be fulfilled, which is collectively described as dataset shift \citep{quionero2009dataset}.

Let $P_{train}(\mathbf{x}, y)$ and $P_{test}(\mathbf{x}, y)$ denote the joint distributions of the training data and the test data. Then, a dataset shift is present when $P_{train}(\mathbf{x}, y) \neq P_{test}(\mathbf{x}, y)$. This discrepancy can stem from changes in the marginal distribution of the input variable $P(\mathbf{x})$, conditional distribution of the target variable $P(y|\mathbf{x})$, or both simultaneously \citep{quionero2009dataset}. Understanding the nature of these dataset shifts is crucial for designing experiments with the goal of evaluating models based on robustness, especially for this work, where we design synthetic data generating processes including shifts in distribution.

\paragraph{Covariate Shift}

Covariate shift describes a situation where the marginal distribution of the input variables differs between training and test set, whereas the conditional distribution of the target remains unchanged \citep{shimodaira2000improving}. We can formalize this as:

\begin{equation}
    P_{tr}(\mathbf{x}) \neq P_{te}(\mathbf{x}) \quad \text{and} \quad P_{tr}(y|\mathbf{x}) = P_{te}(y|\mathbf{x})
\end{equation}

Even if the true target function $f^*$ itself is not changing, a covariate shift can lead to performance degradation, if the hypothesis space $\mathcal{F}$ of a model does not contain the target function $f^*$, i.e. the model is misspecified \citep{wen2014robust}.

In this thesis, we simulate a particular form of covariate shift, here labelled as \textit{noise drift}. It describes a situation, where only non-informative features experience a shift in distribution. This is a purposeful robustness test for a Componentwise Gradient Boosting model, aimed at testing the feature selection mechanism of the algorithm. A robust model must have the ability to ignore irrelevant features to sustain performance loss during a noise drift. Before introducing the Componentwise Gradient Boosting algorithm and its features selection properties in the next section, we formalize the phenomenon of a concept drift more precisely.

\paragraph{Concept Drift}

Concept drift describes a scenario where the conditional distribution of the target variable $P(y|\mathbf{x})$, i.e. the functional relationship, changes between drawing training and test samples \citep{widmer1996learning, gama2014survey}. We formalize this as: 

\begin{equation}
    P_{train}(y|\mathbf{x}) \neq P_{test}(y|\mathbf{x}) \quad \text{and} \quad P_{train}(\mathbf{x}) = P_{test}(\mathbf{x})
\end{equation}

This involves a change of the mapping function $f^*$ itself, trivially leading to performance loss of a model trained on the historical data $\mathcal{D}_{train}$, even if the input distribution $P(\mathbf{x})$ remains unchanged \citep{lu2020learning}. 

In this work, we refer to the concept drift as \textit{meaningful drift}, highlighting the change in the systematic relationship. We simulate this scenario by changing various coefficients of the target functions between training and testing (section 3.1.1.).

\subsection{Component-wise Gradient Boosting (CWB)}

Gradient Boosting Machines, introduced by \citet{friedman2001greedy}, are machine learning algorithms, which construct an additive model using a forward stagewise pattern. Typically, they use deep decision trees (CART) to fit the training data. In contrast, Componentwise Gradient Boosting (CWB) utilizes many simple one-dimensional learners, while associating each of them with just one input feature \citep{buhlmann2003boosting}. 

This constraint makes them highly interpretable, opposed to standard boosting algorithms, which essentially form a black-box predictor. By iteratively updating only one feature component and adding it to the model function, CWB operates directly in the function space and is able to estimate generalized additive models (GAMs) even on high-dimensional data ($p \gg n$) \citep{buhlmann2007boosting}.

\subsubsection{Functional Gradient Descent}

Functional Gradient Descent, the main conceptual framework behind CWB, is a method where the function $\hat{f}(\cdot)$ is directly optimized inside the function space $\mathcal{F}$. This is the core difference to classic Maximum Likelihood Estimation, where function parameters $\boldsymbol{\theta}$ are optimized in a parameter space \citep{mason1999boosting}.

Let the function $\hat{f}$ be the minimizing argument of the empirical risk over the sample data: 

\begin{equation}
    \hat{f} = \underset{f \in \mathcal{F}}{\arg \min} \sum_{i=1}^n L(y_i, f(\mathbf{x}_i)),
\end{equation}

where $L$ is a differentiable, convex function. In this thesis, we aim to find the minimizing function $\hat{f}$ by utilizing the Squared Error Loss ($L_2$ loss) $L(y, f) = \frac{1}{2}(y-f)^2$ \citep{vapnik1995nature}.

The optimization of the model function is conducted iteratively in a additive manner. Let $\hat{f}^{[m-1]}$ be the function estimate at iteration $m-1$. Then, we aim to improve the estimate by adding a function $h(\mathbf{x})$:

\begin{equation}
    \hat{f}^{[m]}(\mathbf{x}) = \hat{f}^{[m-1]}(\mathbf{x}) + \nu \cdot h(\mathbf{x})
\end{equation}
where $0 < \nu \leq 1$ is the step size, known as learning rate. 

Let the \textit{pseudo-residual} $u_i$ at iteration $m$ be given as the negative gradient of the loss function for observation $i$ at iteration $m-1$:

\begin{equation}
    u_i^{[m]} = - \left[ \frac{\partial L(y_i, f(\mathbf{x}_i))}{\partial f(\mathbf{x}_i)} \right]_{f = \hat{f}^{[m-1]}}
\end{equation}

In context of this thesis, we can simplify this to the current residual assuming utilization of the squared error loss:

\begin{equation}
    u_i^{[m]} = - \left( -(y_i - \hat{f}^{[m-1]}(\mathbf{x}_i)) \right) = y_i - \hat{f}^{[m-1]}(\mathbf{x}_i)
\end{equation}

We cannot directly generalize this to new data, since the gradient is only defined on the observations $\mathbf{x}_i$. Hence, we first fit a base learner $h(\mathbf{x}; \boldsymbol{\theta})$ to the pseudo-residuals by using the \textit{least squares criterion}. 


Since the negative gradient is only defined at the observed data points $\mathbf{x}_i$, we cannot generalize to new data directly. Therefore, we fit a parametric base learner $h(\mathbf{x}; \boldsymbol{\beta})$ to the pseudo-residuals using the least squares criterion \citep{friedman2001greedy}:

\begin{equation}
    \hat{\boldsymbol{\theta}}^{[m]} = \underset{\boldsymbol{\beta}}{\arg \min} \sum_{i=1}^n \left( u_i^{[m]} - h(\mathbf{x}_i; \boldsymbol{\theta}) \right)^2
\end{equation}

This results in an updated base learner function $\hat{h}^{[m]}(\mathbf{x}) = h(\mathbf{x}; \hat{\boldsymbol{\theta}}^{[m]})$. The CWB approach restricts this projection further to only one input feature, as presented in the following section.

\subsubsection{The Component-wise Update Step}

Componentwise Gradient Boosting introduces a particular base learner structure, which simplifies them strictly. Let $\mathbf{x} = (x^{(1)}, \dots, x^{(p)})^T$ be the vector of input features. Let $\mathcal{H} = \{h_1, \dots, h_p\}$ be a set of base learners, where each base learner $h_j$ is only defined on the $j$-th input feature $x^{(j)}$ \citep{schmid2008boosting}.

In every iteration, each of the $p$ base learners is fitted to the pseudo-residuals vector $\mathbf{u}^{[m]} = (u_1^{[m]}, \dots, u_n^{[m]})^T$. Then for each feature $j \in \{1, \dots, p\}$, we retain a fitted candidate base learner $\hat{h}_j^{[m]}$ by minimizing the sum of the squared errors:

\begin{equation}
    \hat{h}_j^{[m]} = \underset{h_j \in \mathcal{H}_j}{\arg \min} \sum_{i=1}^n \left( u_i^{[m]} - h_j(x_i^{(j)}) \right)^2
\end{equation}

The standard CWB algorithm then performs a greedy selection of the best base learner by selecting the component $j^*$ which produces the lowest residual sum of squares on the gradient vector: 

\begin{equation}
    j^* = \underset{j \in \{1, \dots, p\}}{\arg \min} \sum_{i=1}^n \left( u_i^{[m]} - \hat{h}_j^{[m]}(x_i^{(j)}) \right)^2
\end{equation}

The update step for iteration $m$ is then following from solely updating the selected component $j^*$:

\begin{equation}
    \hat{f}^{[m]}(\mathbf{x}) = \hat{f}^{[m-1]}(\mathbf{x}) + \nu \cdot \hat{h}_{j^*}^{[m]}(x^{(j^*)})
\end{equation}

The whole procedure is summarized in Algorithm \ref{alg:cwb}. We finally observe the additive structure of the final model, which is produced by accumulation of the incremental updates for each feature: 

\begin{equation}
\hat{f}(\mathbf{x}) = \hat{\beta}_0 + \sum_{j=1}^p \hat{f}_j(x^{(j)})
\end{equation}

where each component $\hat{f}_j$ aggregates all boosting updates for the $j$-th input feature \citep{hastie2009elements}.

\begin{algorithm}
\caption{Component-wise Gradient Boosting with $L_2$ Loss}
\label{alg:cwb}
\begin{algorithmic}[1]
\State \textbf{Initialize:} $\hat{f}^{[0]}(\mathbf{x}) \equiv \bar{y} = \frac{1}{n}\sum_{i=1}^n y_i$ (or intercept 0).
\State \textbf{Set:} Learning rate $0 < \nu \le 1$, number of iterations $M_{stop}$.
\For{$m = 1$ to $M_{stop}$}
    \State \textbf{Compute pseudo-residuals:} 
    \State \quad $u_i^{[m]} = y_i - \hat{f}^{[m-1]}(\mathbf{x}_i) \quad \text{for } i=1,\dots,n$
    \State \textbf{Fit base learners:}
    \For{$j = 1$ to $p$}
        \State Fit learner $h_j$ to $\mathbf{u}^{[m]}$ to obtain estimate $\hat{h}_j$.
        \State Compute loss $S_j = \sum_{i=1}^n (u_i^{[m]} - \hat{h}_j(x_i^{(j)}))^2$.
    \EndFor
    \State \textbf{Select best component:}
    \State \quad $j^* = \arg \min_{j} S_j$
    \State \textbf{Update model:}
    \State \quad $\hat{f}^{[m]}(\mathbf{x}) = \hat{f}^{[m-1]}(\mathbf{x}) + \nu \cdot \hat{h}_{j^*}(x^{(j^*)})$
\EndFor
\State \textbf{Return} $\hat{f}^{[M_{stop}]}(\mathbf{x})$
\end{algorithmic}
\end{algorithm}

\subsubsection{Implicit Feature Selection Properties}

The inherent feature selection capabilities of the component-wise update procedure are a critical advantage of the CWB algorithm over other boosting algorithms. In each iteration it only selects the most informative feature $j^*$ potentially leading to some non-informative features being left out entirely from the final model \citep{hofner2014model}.

Hence, when regularizing the algorithm with early stopping, the final model is relatively sparse. We denote the set of indices features, which were selected at least once during training, with $\mathcal{J}_{selected}$. Then, the estimated functions of unselected features stay zero:

\begin{equation}
    \forall j \notin \mathcal{J}_{selected}: \quad \hat{f}_j(x^{(j)}) = 0
\end{equation}

Shown by \citet{efron2004least}, this characteristic of the CWB algorithm produces an close link to Lasso ($L_1$-regularized) regression, which states that for infinitesimal learning rates ($\nu \to 0$), the parameter paths of component-wise linear boosting approximate the regularization path of Lasso regression. However, while Lasso uses an explicit penalty term, CWB obtains the regularization implicitly via the number of iterations $M_{stop}$. Hence, CWB becomes particularly suitable for high-dimensional scenarios ($p \gg n$) \citep{buhlmann2002boosting}, such as in the Riboflavin dataset employed in this thesis.

\subsection{Base Learners}
\label{sec:base_learners}

After establishing the properties of the selection process of base learners, we proceed by formalizing the different base learner types employed in this work, while stating their characteristics. We distinguish between two categories of base learners: parametric learners (Linear, Polynomial) and non-parametric learners (B-Splines, Decision Trees)

\subsubsection{Parametric Learners}
\label{sec:parametric}


Parametric base learners assume the relationship between the input feature $x^{(j)}$ and pseudo-residuals $u^{[m]}$ to follow a function, which can be captured using a finite set of parameters \citep{hastie2009elements}. They are particularly useful to capture global relationships and smooth target functions, while offering a high interpretability and lower variance compared to non-parametric base learners \citep{buhlmann2003boosting}.

\paragraph{Linear Base Learners}

The linear base learner marks the most fundamental base learner in CWB. Combined with the squared error loss, it is often referred to as $L_2$-Boosting \citep{buhlmann2007boosting}. For a feature $j$, let $h_j(x^{(j)}; \theta)$ be the linear base learner, which approximates the negative gradient vector with a simple linear function:
\begin{equation}
    h_j(x^{(j)}; \theta) = \theta \cdot x^{(j)}
\end{equation}
where $\theta \in \mathbb{R}$ is a scalar slope parameter. The global intercept parameter $\hat{\theta}_0$ is handled separately, by initializing it in the first step of the CWB algorithm (Algorithm \ref{alg:cwb}), leading to linear learners optimizing only the slope parameter \citep{hofner2014model}. The minimization of the square error loss is conducted by computing the Ordinary Least Squares (OSL) solution: 
\begin{equation}
    \hat{\theta} = \frac{(\mathbf{x}^{(j)})^T \mathbf{u}^{[m]}}{(\mathbf{x}^{(j)})^T \mathbf{x}^{(j)}} = \frac{\sum_{i=1}^n x_i^{(j)} u_i^{[m]}}{\sum_{i=1}^n (x_i^{(j)})^2}
\end{equation}
where $\mathbf{x}^{(j)}$ is the observation vector for the $j$-th input feature and $\mathbf{u}^{[m]}$ denotes the pseudo-residuals vector at iteration $m$ \citep{buhlmann2007boosting}.


\paragraph{Polynomial Base Learners}

For non-linear relationships and curvature in the data, which can not be modelled with linear term, we utilize polynomial base learners. They extend the hypothesis space via basis expansion of the input feature $x^{(j)}$ \citep{hastie2009elements}.

Let the polynomial base learner $h_j(x^{(j)}; \boldsymbol{\theta})$, for feature $j$ and a degree $d$ be described as linear combination of polynomial power functions:

\begin{equation}
    h_j(x^{(j)}; \boldsymbol{\theta}) = \theta_0 + \sum_{k=1}^d \theta_k (x^{(j)})^k
\end{equation}

where $\boldsymbol{\theta} = (\theta_0, \theta_1, \dots, \theta_d)^T \in \mathbb{R}^{d+1}$ denotes the coefficient vector. 

To obtain the optimal parameter vector $\hat{\boldsymbol{\theta}}$, we minimize the sum of squared errors between polynomial approximation and pseudo-residuals. This is done with the normal equations for the expanded design matrix denoted by $\mathbf{X}_j = [\mathbf{1}, \mathbf{x}^{(j)}, (\mathbf{x}^{(j)})^2, \dots, (\mathbf{x}^{(j)})^d]$:

\begin{equation}
    \hat{\boldsymbol{\theta}} = (\mathbf{X}_j^T \mathbf{X}_j + \lambda \mathbf{I})^{-1} \mathbf{X}_j^T \mathbf{u}^{[m]}
\end{equation}

where $\mathbf{I}$ is the identity matrix and $\lambda$ is a small regularization term (e.g., $\lambda = 10^{-8}$) which ensures numerical stability for the inversion of $\mathbf{X}_j^T \mathbf{X}_j$ \citep{hofner2014model}. 

Polynomial base learners offer more flexibility than linear base learners, while both are parametric. Nevertheless, their global nature also implies that updates in one region of the feature space $x^{(j)}$ lead to changes of the function estimate across the whole feature space. This distinguishes both the linear and polynomial base learner from local methods such as B-Splines and decision trees \citep{buhlmann2003boosting}.

\subsubsection{Non-Parametric Learners}
\label{sec:non_parametric}

In contrast to parametric base learners, which in the linear and polynomial case are globally defined, non-parametric learners do not rely on a finite number of parameters or a global shape. This makes them flexible enough to capture local structure in the data, a scenario where global learners often fail. They are able to grow in complexity with more data available \citep{hastie2009elements}.

In CWB, non-parametric learners are fitted to specific regions of the input feature space $x^{(j)}$. Especially, in case of sharp changes and high-frequent patterns, they are able to adapt to the data, where global learners lack flexibility. However, this desired flexibility comes with the high price of overfitting risk and variance. Therefore, they require robust regularization, such as explicit penalization \citep{hofner2014model}, which is discussed in section 2.4. In this thesis, we apply two different types of non-parametric learners: B-Splines and Decision Stumps.

\paragraph{B-Spline Basis Functions}

A B-Spline (Basis Spline) is a piecewise polynomial function, which is defined on a set of intervals. At the interval boundaries, also known as \textit{knots}, these polynomial pieces are connected smoothly \citep{Boor1986BasicSplineB}.

Let $\boldsymbol{\xi} = (\xi_1, \xi_2, \dots, \xi_K)$ denote a sequence of knots that divide the space of $x^{(j)}$. In this work, we place the knots according to fixed quantiles of the input feature distribution. This method ensures that the data is evenly covered over the intervals, while reducing variance of the the fitting process in sparse data regions. 

Let the B-Spline learner $h_j(x^{(j)}; \boldsymbol{\theta})$ with degree $d$ be defined as linear combination of $M = K + d - 1$ ($K$ knots) basis functions $B_{k,d}(x^{(j)})$:

\begin{equation}
    h_j(x^{(j)}; \boldsymbol{\theta}) = \sum_{k=1}^{M} \theta_k B_{k,d}(x^{(j)})
\end{equation}

where $\boldsymbol{\theta} = (\theta_1, \dots, \theta_M)^T$ denotes the vector of coefficients. The basis functions $B_{k,d}$ can be constructed recursively with the Cox-de Boor formula. The construction of a basis function of $d=k$ starts from the defined basis function for $d=0$:

For degree $d=0$ (piecewise constants), the basis functions are defined as:

\begin{equation}
    B_{k,0}(x) = 
    \begin{cases} 
    1 & \text{if } \xi_k \leq x < \xi_{k+1} \\
    0 & \text{otherwise}
    \end{cases}
\end{equation}

Then, for $d > 0$, the basis functions are recursively computed as \citep{zhou2023tutorial}:

\begin{equation}
    B_{k,d}(x) = \frac{x - \xi_k}{\xi_{k+d} - \xi_k} B_{k,d-1}(x) + \frac{\xi_{k+d+1} - x}{\xi_{k+d+1} - \xi_{k+1}} B_{k+1,d-1}(x)
\end{equation}

In this work, we employ Cubic B-Splines ($d=3$), which guarantees second order continuity at the knots, leading to a smooth approximation.

Similar to the learners discussed in Section \ref{sec:parametric}, we estimate the coefficients $\boldsymbol{\theta}$ by minimizing the sum of squared errors against the pseudo-residuals $u^{[m]}$. We denote the $n \times M$ design matrix with $\mathbf{B}_j$ where the $(i, k)$-th entry is $B_{k,3}(x_i^{(j)})$. The optimal coefficient vector is derived from the least squares solution \citep{hofner2014model}: 

\begin{equation}
    \hat{\boldsymbol{\theta}} = (\mathbf{B}_j^T \mathbf{B}_j)^{-1} \mathbf{B}_j^T \mathbf{u}^{[m]}
\end{equation}

\paragraph{Decision Stumps}

Decision stumps are the simplest special form of decision trees, with a restricted depth of $d=1$. They approximate local constants by dividing the feature space $x^{(j)}$ into two disjoint regions. Let  $h_j(x^{(j)}; \boldsymbol{\theta})$ be a decision stump for the $j$-th feature \citep{breiman1984cart}. It can be formalized as piecewise constant function:

\begin{equation}
    h_j(x^{(j)}; \boldsymbol{\theta}) = c_L \cdot \mathbb{I}(x^{(j)} \leq t) + c_R \cdot \mathbb{I}(x^{(j)} > t)
\end{equation}

where $t$ is the splitting threshold, and $\boldsymbol{\theta} = \{t, c_L, c_R\}$ denotes the set of parameters. The constants $c_L$ and $c_R$ denote the estimated values for the respective region. With the utilisation of the $L_2$ loss, the optimal values for $c_L$ and $c_R$ for a threshold $t$ are simply computed by the average of the pseudo-residuals $u_i^{[m]}$ of the respective region \citep{friedman2001greedy}:

\begin{equation}
c_L = \frac{\sum_{i: x_i^{(j)} \leq t} u_i^{[m]}}{\sum_{i: x_i^{(j)} \leq t} 1}, \quad
c_R = \frac{\sum_{i: x_i^{(j)} > t} u_i^{[m]}}{\sum_{i: x_i^{(j)} > t} 1}
\end{equation}

A single decision stump holds high bias and is not able to capture complex patterns, but including them in a CWB ensemble enables the model to approximate non-linear functions or discontinuities, such as step functions, a scenario where, e.g. global polynomials struggle \citep{hastie2009elements}.

\subparagraph{Histogram-based Optimization (Binning)}

To find the optimal split $t$ which reduces the $L_2$ loss, the standard CWB algorithm requires sorting all feature values $x^{(j)}$ at each iteration. This comes with a computational complexity of $\mathcal{O}(n \log n)$ per feature \citep{schalk2022accelerated}. In order to improve computational efficiency and scalability in high-dimensional settings, we utilize a histogram-based approximation method, which discretizes the continuous input feature space into $B$ bins. This reduces the candidates for $t$ to $B-1$ \citep{ke2017lightgbm}.

Let $\Phi: \mathbb{R} \to \{1, \dots, B\}$ be a mapping function which assigns each observation $x_i^{(j)}$ to a bin index $b_i = \Phi(x_i^{(j)})$ based on pre-computed quantiles of the input feature distribution. We then aggregate the pseudo-residuals $u_i^{[m]}$ into the histogram bins and compute the following statistics for each bin $k \in \{1, \dots, B\}$:

\begin{equation}
    G_k = \sum_{i: b_i = k} u_i^{[m]}, \quad N_k = \sum_{i: b_i = k} 1
\end{equation}

where $G_k$ denotes the sum of pseudo-residuals and $N_k$ the sample count in the $k$-th bin. We now can evaluate the quality of the split candidates $t$ with a computational complexity of $\mathcal{O}(B)$, while being independent of the number of samples $n$ \citep{chen2016xgboost}. Therefore, cumulative sums are computed over the quantiles of the histogram. Then the optimal split bin $k^*$ is set to maximize the reduction of the sum of squared errors:

\begin{equation}
    Gain(k) = \frac{(\sum_{l=1}^k G_l)^2}{\sum_{l=1}^k N_l} + \frac{(\sum_{r=k+1}^B G_r)^2}{\sum_{r=k+1}^B N_r}
\end{equation}

This discretization technique, primarily introduced to accelerate the fitting process and reduce memory consumption \citep{schalk2023accelerated}, also produces implicit regularization, which is achieved by limiting the number of candidate split points, thus reducing the variance of the base learner \citep{ke2017lightgbm}.


\subsection{Competing Base Learners and Model Complexity}

The base learners introduced in Section \ref{sec:base_learners} offer strong tools to approximate particular shapes of target functions. Nevertheless, in standard CWB implementations, the users it typically required to specify only one base learner type for all input features. In many high-dimensional real-world scenarios however, the underlying relationships between input and target variable cannot be captured by single functional forms, because of their heterogeneous nature. While some features may have a linear influence on the target, others may follow complex non-linear or discontinued patterns. This highlights a trade-off, which can be solved by expanding the hypothesis space of the CWB algorithm \citep{schmid2008boosting}. Instead of fitting just one base learner type to each input feature ${x^{(j)}}$ at each iteration $m$, we extend the space to a set of base learner types $\mathcal{T} = \{\text{linear}, \text{polynomial}, \text{spline}, \text{tree}\}$ for each feature ${x^{(j)}}$. By then fitting a base learner of each type, the optimization at iteration $m$ becomes a joint optimization problem over the feature index $j$ and the learner type $\tau \in \mathcal{T}$:

\begin{equation}
    (j^*, \tau^*) = \underset{j \in \{1, \dots, p\}, \tau \in \mathcal{T}}{\arg \min} \sum_{i=1}^n \left( u_i^{[m]} - \hat{h}_{j, \tau}(x_i^{(j)}) \right)^2
\end{equation}

Due to the differing complexity levels of the learner types, this diversification enables the CWB algorithm to determine an appropriate complexity for each input feature ${x^{(j)}}$. However, this flexibility thereby introduces a structural bias in the feature selection problem, which must be addressed to prevent the model from overfitting \citep{hofner2014model}.

\subsubsection{The Biased Selection Problem}

The problem of introducing bias when enabling base learners of varying flexibility and complexity to compete with each other is known as \textit{variable selection bias}. Let $\hat{h}_{j, \text{linear}}$ be a linear base learner and $\hat{h}_{j, \text{spline}}$ be a B-Spline learner for the feature $j$. The flexibility of the learners can be described with their degrees of freedom $df$, which can be derived from the trace of their projection matrix, that maps the observations to the fitted values \citep{hofner2011selection}. While the linear base learner has $df_{lin} = 1$ (excluding the intercept parameter), the cubic B-Spline has $df_{spline} = K + 3$ where $K$ is the number of knots. The hypothesis space of the linear base learner is a subspace of the hypothesis space of the B-Spline learner, which makes them fundamentally more flexible, and allowing them to fit a closer approximation of the random noise in $\mathbf{u}^{[m]}$. Therefore, the more flexible base learner will often produce a smaller residual sum of squares, even if the underlying relationship of $x^{(j)}$ and $y$ is strictly linear \citep{hastie2009elements}:

\begin{equation}
    \sum_{i=1}^n \left( u_i^{[m]} - \hat{h}_{j, \text{spline}}(x_i^{(j)}) \right)^2 \leq \sum_{i=1}^n \left( u_i^{[m]} - \hat{h}_{j, \text{linear}}(x_i^{(j)}) \right)^2
\end{equation}

Hence, the greedy CWB feature selection approach will then systematically prefer the more complex base learners. If the algorithm relies on this residual sum of squares solely, two unfavourable effects are induced:

\begin{enumerate}
    \item \textbf{Overfitting:} The model fits to the random noise, thus increasing the variance of the predictive performance. 
    \item \textbf{Lower interpretability:} Because of the misspecification of the underlying relationships, the model becomes unnecessarily complex, and looses its inherent interpretability properties.
\end{enumerate}

Therefore, we adjust the feature selection algorithm to allow fair competition between the base learners. \citet{hofner2011selection} proposed that this requires to constraint the degrees of freedom of the different learners to a comparable level or to use specific penalization methods. In this thesis, we proceed with this by utilizing orthogonal decomposition and Ridge penalization, as discussed in the following sections.

\subsubsection{Orthogonal Decomposition for Identifiability}
\label{sec:orth_decomp}

As previously stated, one factor leading to variable selection bias is the shared functional hypothesis subspace of competing base learners. A complex non-linear base learner, such as the B-Splines, is perfectly capable of reducing itself to a linear function. Simultaneously, the flexible base learner is able to capture not only a linear signal, but also a non-linear deviation plus noise. The linear base learner becomes statistically obsolete in such settings, while also becoming unidentifiable \citep{hofner2011selection}.

Therefore, we employ an approach, which decomposes the function space of a complex learner into a subspace which is spanned by the linear learner and a orthogonal subspace, referring to the non-linear part. Let $\mathbf{X}_j = [\mathbf{1}, \mathbf{x}^{(j)}]$ denote the design matrix of a linear learner

Let $\mathbf{X}_j = [\mathbf{1}, \mathbf{x}^{(j)}]$ denote the design matrix of a linear base learner for feature $j$, including the intercept column. Note that, contrary to our formalization of the standard CWB algorithm, here the intercept is explicitly included in the decomposition step. This serves the reason to ensure that the projection captures the whole linear subspace, therewith leaving the orthogonal complementary with the strictly non-linear part. Let $\mathbf{B}_j$ be the design matrix of the complex base learner (e.g. B-Spline basis matrix in Section \ref{sec:non_parametric}). The flexible design matrix is projected onto the linear subspace using the projection matrix $\mathbf{P}_j$:

\begin{equation}
    \mathbf{P}_j = \mathbf{X}_j (\mathbf{X}_j^T \mathbf{X}_j)^{-1} \mathbf{X}_j^T
\end{equation}

The original design matrix $\mathbf{B}_j$ then can be decomposed into a projected part $\mathbf{B}_j^{proj} = \mathbf{P}_j \mathbf{B}_j$ and an orthogonal part $\tilde{\mathbf{B}}_j$:

\begin{equation}
    \tilde{\mathbf{B}}_j = (\mathbf{I} - \mathbf{P}_j) \mathbf{B}_j = \mathbf{B}_j - \mathbf{X}_j \underbrace{(\mathbf{X}_j^T \mathbf{X}_j)^{-1} \mathbf{X}_j^T \mathbf{B}_j}_{\boldsymbol{\Gamma}_j}
\end{equation}

Here, $\boldsymbol{\Gamma}_j$ denotes the coefficients for the projection of the non-linear design matrix onto the linear design matrix. By using the decomposed design matrix $\tilde{\mathbf{B}}_j$ for the non-linear learner, we ensure that the columns of $\tilde{\mathbf{B}}_j$ are orthogonal to the columns of $\mathbf{X}_j$, and thus, restricting the non-linear learner for modelling only non-linear patterns \citep{schmid2008boosting}.

This leads to a clear additive representation of the approximated feature effect:

\begin{equation}
    \hat{\mathbf{f}}_j = \underbrace{\mathbf{X}_j \hat{\boldsymbol{\theta}}_{lin}}_{\text{linear part}} + \underbrace{\tilde{\mathbf{B}}_j \hat{\boldsymbol{\theta}}_{nonlin}}_{\text{non-linear deviation}}
\end{equation}

This shows, that identifiability is restored, utilizing orthogonal decomposition for separation of the hypothesis spaces of the learner. Now, for the case of a linear signal, the orthogonal non-linear part $\tilde{\mathbf{B}}_j$ will only capture non-linear deviations and the linear part $\mathbf{X}_j$ will capture the signal. However, orthogonality of hypothesis spaces alone does not eliminate the feature selection bias emerging from the high-dimensionality of $\tilde{\mathbf{B}}_j$. This higher flexibility leaves the non-linear base learner with an competitive advantage, which must be constrained to ensure a fair feature selection. 


\subsubsection{Degrees of Freedom and Ridge Penalization}
\label{sec:df_ridge}

Even after applying orthogonal decomposition, the non-linear design matrix $\tilde{\mathbf{B}}_j$ usually has a higher rank than the linear design matrix $\mathbf{X}_j$, allowing higher flexibility when fitting noise. To assure a fair feature selection, we employ Ridge penalization ($L2$-regularization) to the non-linear learners, explicitly regulating their effective degrees of freedom $df$ \citep{hofner2011selection}.

Therefore, the least squares criterion is modified by adding a regularizing penalty term:

\begin{equation}
    \hat{\boldsymbol{\theta}} = \underset{\boldsymbol{\theta}}{\arg \min} \sum_{i=1}^n \left( u_i^{[m]} - \tilde{\mathbf{B}}_j \boldsymbol{\theta} \right)^2 + \lambda_j \boldsymbol{\theta}^T \boldsymbol{\Omega} \boldsymbol{\theta}
\end{equation}

where $\lambda_j \geq 0$ denotes a regularization parameter and $\boldsymbol{\Omega}$ denotes the penalty matrix. The shape and structure of $\boldsymbol{\Omega}$ depends on the base learner type. For polynomials, we use a standard Ridge penalization where $\boldsymbol{\Omega} = \mathbf{I}$. The penalty term then penalizes the square Euclidean norm of the learner coefficients, leading to shrinkage of the coefficients: 

\begin{equation}
    \lambda_j \boldsymbol{\theta}^T \mathbf{I} \boldsymbol{\theta} = \lambda_j \sum_{k} \theta_k^2
\end{equation}

For B-Spline base learners, we use P-Splines (Penalized B-Splines) as introduced by \citep{eilers1996flexible}. Let $\mathbf{D}_d$ be the $d$-th order difference matrices, then we construct $\boldsymbol{\Omega}$ by $\boldsymbol{\Omega} = \mathbf{D}_d^T \mathbf{D}_d$. This penalizes the squared deviation of adjacent B-Spline coefficients, thereby regularizing towards smoothness of the approximated function. The penalty term is derived by:

\begin{equation}
    \lambda_j \boldsymbol{\theta}^T \mathbf{D}_d^T \mathbf{D}_d \boldsymbol{\theta} = \lambda_j ||\mathbf{D}_d \boldsymbol{\theta}||^2 = \lambda_j \sum_{k} (\Delta^d \theta_k)^2
\end{equation}

In this thesis, we use second-order differences ($d=2$) to penalize all deviations from linearity. Using larger values for $d$, allows the Spline base learner for more flexibility. Now, the general least squares solution for the penalized coefficients becomes:

\begin{equation}
    \hat{\boldsymbol{\theta}} = (\tilde{\mathbf{B}}_j^T \tilde{\mathbf{B}}_j + \lambda_j \boldsymbol{\Omega})^{-1} \tilde{\mathbf{B}}_j^T \mathbf{u}^{[m]}
\end{equation}

While we established a penalized criterion which yields estimates for the base learner coefficients, the regularization parameter $\lambda_j$ lacks a determining framework. Therefore we use the concept of effective degrees of freedom. Let $\hat{\mathbf{y}} = \mathbf{S}_{\lambda_j} \mathbf{u}^{[m]}$ describe the fitted values for a base learner $j$, where $\mathbf{S}_{\lambda_j}$ is the penalized projection matrix: 

\begin{equation}
    \mathbf{S}_{\lambda_j} = \tilde{\mathbf{B}}_j (\tilde{\mathbf{B}}_j^T \tilde{\mathbf{B}}_j + \lambda_j \boldsymbol{\Omega})^{-1} \tilde{\mathbf{B}}_j^T
\end{equation}

Let the effective degrees of freedom $df(\lambda_j)$ be defined as the trace of the penalized projection matrix $\mathbf{S}_{\lambda_j}$:

\begin{equation}
    df(\lambda_j) = \text{tr}(\mathbf{S}_{\lambda_j})
\end{equation}

We see then, that $df(\lambda_j)$ decreases strictly monotonically with $\lambda_j$. With a increasing $\lambda_j$, the penalty term $\lambda_j \boldsymbol{\Omega}$ progressively dominates the inverse matrix, resulting in suppression of the eigenvalues of $\mathbf{S}_{\lambda_j}$, and thus reducing the trace. Hence, we can solve for $\lambda_j^*$ for each desired flexibility level $df_{target}$ with:

\begin{equation}
    df(\lambda_j^*) = df_{target}
\end{equation}

This regularized mechanism allows the CWB Boosting algorithm to distinguish between linear and non-linear features while enforcing a fair variable selection.

\subsection{Regularization Techniques}

In context of boosting, regularization is inevitably needed to prevent the model from overfitting the random noise $\epsilon$ in the training data. Because of its iterative minimization of the empirical risk, unconstrained boosting models will converge to a estimation which interpolates the training observations, and thus maximizing variance and degradation in predictive performance  \citep{buhlmann2007boosting}. To mitigate this limitation, CWB uses both implicit and explicit regularization mechanisms. Implicit regularization is given by algorithm itself in the form of the shrinkage parameter (learning rate $\nu$) and the number of boosting iterations $M_{stop}$ \citep{friedman2001greedy}. Explicit regularization techniques involve the modification of the objective function or other optimization constraints, such as the discussed penalization of base learners (Section \ref{sec:df_ridge}) or Flooding regularization discussed din Section \ref{sec:flooding}. 

\subsubsection{Early Stopping}

Early stopping is the main regularizer in CWB, and follows from the empirical observation, that the generalization error of a U-shaped curve with respect to the iterations $m$. In the first part of the trajectory the error decreases while the model learns the true target function and minimizing bias, then reaching a global minimum, before then increasing as the model starts to fit the random noise and thus increasing the variance \citep{zhang2005boosting}.

Now, the iteration index $m$ is utilized as inverse regularization parameter, which when set to $m=0$ corresponds to a maximum regularization, while $m \to \infty$ corresponds to an unregulated fully converged model \citet{buhlmann2003boosting}. Therefore, deciding for an optimal stopping iteration $m_{stop}$ is subject optimization.

\paragraph{Experimental Implementation}

In this thesis, we employ K-Fold Cross-Validation (Section \ref{sec:cv}) to determine the optimal stopping iteration $m_{stop}$, by applying the K-Fold strategy on the training set in every single training process of the experiment.

Let $K$ be the number of folds. For each fold $k$, we train a model $\hat{f}^{(k)}$ on the respective data partition and compute the validation loss trajectory over the iterations $m = 1, \dots, M_{max}$ on the held-out validation fold. We then aggregate these trajectories to compute the cross-validation error by averaging the empirical risk, i.e. Mean Squared Error (MSE), over all folds in each iteration:

\begin{equation}
    \text{CV}(m) = \frac{1}{K} \sum_{k=1}^K \mathcal{R}_{val}^{(k)}(m)
\end{equation}

The optimal stopping iteration is then defined as the index which minimizes the cross-validation error:
\begin{equation}
    m_{stop} = \underset{m}{\arg \min} \; CV(m)
\end{equation}

The final model is then fitted by using the complete training set (combining all $K$ folds) and iterating to the optimal stopping point $m_{stop}$. Thereby, we reduce the bias towards certain training data splits, while also preventing data leakage from the test set \citep{bischl2012resampling}. 

\subsubsection{Flooding Regularization}
\label{sec:flooding}

The previously discussed regularization technique of early stopping intends to prevent overfitting by terminating the fitting process before the model fits the noise. It does not change the optimization objective itself. 

In contrast, \textit{Flooding} proposed by \citet{ishida2020need} is regularization method, which explicitly modifies the optimization objective with an update to the empirical risk function. It thereby prevents the training loss from dropping below a pre-defined threshold, described as \textit{flood level} $b$.

The theoretical justification for flooding regularization comes from the observation that modern models often have the capacity to achieve a near-zero training loss, which essentially implies memorization of training data. \citet{ishida2020need} argue that incentivizing this training behaviour is unnecessary and potentially even harmful in terms of generalization. By introducing the flood level $b$ the model is forced to stay above a certain non-zero threshold, and convicted to 'float' around this level. It is hypothesized, that the potentially resulting 'random walk' of the model around the flood level, drives it into flatter minima of the loss landscape, which are related to better generalization \citep{hochreiter1997flat}. 

\paragraph{Gradient Definition under Flooding}

Let the flooded empirical risk $\tilde{\mathcal{R}}_{emp}(f)$ be defined as:

\begin{equation}
    \tilde{\mathcal{R}}_{emp}(f) = |\mathcal{R}_{emp}(f) - b| + b
\end{equation}

where $b > 0$ is the flood level \citep{ishida2020need}. During optimization this flooded risk is minimized. The gradient of $\tilde{\mathcal{R}}_{emp}(f)$ with respect to the model function $f$ is derived as:

\begin{equation}
    \nabla \tilde{\mathcal{R}}_{emp}(f) = \text{sign}(\mathcal{R}_{emp}(f) - b) \cdot \nabla \mathcal{R}_{emp}(f)
\end{equation}

This leads to two distinct phases of the optimization \citep{ishida2020need}:
\begin{itemize}
    \item \textbf{Gravity phase ($\mathcal{R}_{emp}(f) > b$):} The gradient points in the same direction as the original risk. The algorithm performs standard gradient descent.
    \item \textbf{Buoyancy phase ($\mathcal{R}_{emp}(f) < b$):} The gradient sign is flipped. The algorithm effectively performs \textit{gradient ascent}, while actively trying to increase the loss to drive it back to $b$.
\end{itemize}

In the context of CWB, this is utilized by adjusting the pseudo-residuals $\mathbf{u}^{[m]}$. If the empirical risk $\mathcal{R}_{emp}$ of the current ensemble $\hat{f}^{[m-1]}$ falls below $b$, the sign of the residuals is inverted:

\begin{equation}
    \mathbf{u}^{[m]} = \begin{cases} 
    -\nabla L(\mathbf{y}, \hat{f}^{[m-1]}(\mathbf{x})) & \text{if } \mathcal{R}_{emp}(\hat{f}^{[m-1]}) > b \\
    +\nabla L(\mathbf{y}, \hat{f}^{[m-1]}(\mathbf{x})) & \text{if } \mathcal{R}_{emp}(\hat{f}^{[m-1]}) < b 
    \end{cases}
\end{equation}

\paragraph{Theoretical Justification}

The theoretical justification for flooding stems from the Mean Squared Error (MSE) of the risk estimator. \citet{ishida2020need} provide a theorem stating that if the flood level $b$ is chosen such that $b \leq R(f^*)$, where $R(f^*)$ is the true risk, i.e. true noise variance, the MSE of the flooded empirical risk estimator is lower or equal to that of the original empirical risk estimator:

\begin{equation}
    \text{MSE}(\tilde{\mathcal{R}}_{emp}(f)) \leq \text{MSE}(\mathcal{R}_{emp}(f))
\end{equation}

It follows that flooding leads to variance reduction for the risk estimator, while preventing the model to overfit to specific realizations of training noise, if the flood level $b$ approximates the true noise variance from below.

\paragraph{Adaptation to Component-wise Boosting}

In their work, \citet{ishida2020need} showed the behaviour of flooding using Stochastic Gradient Descent (SGD). They observe, that the 'random walk' effect around the flood level, is mostly induced by the stochastic noise introduced via the mini-batch sampling method. 

While \citet{friedman2002stochastic} shows that Stochastic Gradient Descent (SGD) in boosting can improve generalization, there may be a conflict for the objective of \textit{Componentwise} Gradient Boosting, which is accurate and interpretable feature selection. When introducing variance into the residuals via mini-batch subsampling, it may be causing the algorithm to select features based on sampling artefacts, and not based on the underlying signal. Therefore, in this work, flooding is adapted in a full-batch setting.

However, the deterministic full-batch approach may limit flooding regularization method to unveil their potential of exploring the hypothesis space by 'randomly walk' around the flood level, and thus potentially escape local minima. Therefore, we introduce a controlled source of stachasticity with the Top-K feature selection discussed in Section \ref{sec:topk}. It has to be distinguished that SGD adds noise into the \textit{estimation} of feature quality, whereas Top-K feature selection adds variance into the \textit{selection} among already high-ranked candidates.



\subsection{Advanced Feature Selection Mechanisms}
\label{sec:adv_selection}

The greedy feature selection mechanism in CWB offers intrinsic variable selection as well as high interpretability of the final model. However, in high-dimensional scenarios with the presence of high noise and correlations among features, it can be a victim to instability \citep{Hofner2015}. We aim at mitigating these limitations by introducing two advanced feature selection mechanisms: \textit{Top-K Stochastic Selection} and \textit{Momentum-based Selection}.

\subsubsection{Limitations of Greedy Selection}

The standard CWB algorithm selects the best feature $j^*$ in iteration $m$ by minimizing the residual sum of squares:
\begin{equation}
    j^* = \underset{j \in \{1, \dots, p\}}{\arg \min} \sum_{i=1}^n \left( u_i^{[m]} - \hat{h}_j(x_i^{(j)}) \right)^2
\end{equation}

This approach can lead to problems in scenarios where several features have similar predictive power, which is usually due to correlation of their marginal distributions $P(x^{(j)})$. This can lead to the algorithm switching between correlated features over iterations \citep{Hofner2015}. Furthermore, by strictly following the steepest descent in the hypothesis space, like greedy selection does, the algorithm may ignore steps in early iterations, which seem suboptimal, but lead to a better final model at the end \citep{efron2004least, hastie2009elements}.

\subsubsection{Top-K Stochastic Selection}
\label{sec:topk}

To induce a explorative behaviour into the feature selection process, and thus into the optimization trajectory, we implement a stochastic selection mechanism, referred to as Top-K selection. Here, the we move away from deterministic feature selection, and introduce a method which selects the feature from a bucket of $K$ best candidate features.

Let $\mathcal{S}^{[m]} = \{S_1^{[m]}, \dots, S_p^{[m]}\}$ be the set of squared error losses for all features in iteration $m$, where $S_j^{[m]} = \sum_{i=1}^n (u_i^{[m]} - \hat{h}_j(x_i^{(j)}))^2$. Then, the set of indices of the $K$ best features $\mathcal{J}_K$ is defined, such that:

\begin{equation}
    \forall j \in \mathcal{J}_K, \forall l \notin \mathcal{J}_K: \quad S_j^{[m]} \le S_l^{[m]} \quad \text{and} \quad |\mathcal{J}_K| = K
\end{equation}

The selected feature $j^*$ is drawn from $\mathcal{J}_K$ with a probability distribution, depending on the rank of the loss of the feature. In our implementation, we use a rank-weighted mass function with the objective to favour higher ranked features, but still conceding exploration. Let $r_j \in \{1, \dots, K\}$ be the rank of the feature $j$, where 1 corresponds to the lowest loss, i.e. the best feature. Then the probability of selecting the feature $j$ is given by:
\begin{equation}
    P(j^* = j | j \in \mathcal{J}_K) = \frac{w_{r_j}}{\sum_{k=1}^K w_k}, \quad \text{with } w_r = K - r + 1
\end{equation}
Setting $K=1$, this method becomes equal to the standard greedy algorithm of CWB. For $K > 1$, it introduces controlled stochasticity to the selection mechanism.

\subsubsection{Momentum-based Feature Selection}

The standard CWB algorithm is memoryless in terms of feature sections history, which means the decision at iteration $m$ is only depending on the current residuals $\mathbf{u}^{[m]}$. With Momentum-based selection we introduce a history term, which is supposed to smooth the selection process. 

We establish a momentum vector $\mathbf{v} \in \mathbb{R}^p$ initialized to zero. At each iteration, we define the gain $g_j^{[m]}$ for each feature $j$, which quantifies its provided improvement. For numerical stability across different residual scales, we define it relative to the feature with the lowest improvement and normalize:

\begin{equation}
    g_j^{[m]} = \frac{\max_k (S_k^{[m]}) - S_j^{[m]}}{\max_k (\max_l (S_l^{[m]}) - S_k^{[m]}) + \epsilon}
\end{equation}

with $k,l \in \set{1, \dots ,p}$. The current momentum $v_j^{[m]}$ of feature $j$ is updated at each iteration using an exponential moving average with a decay rate $\alpha \in [0, 1)$:

\begin{equation}
    v_j^{[m]} = \alpha \cdot v_j^{[m-1]} + g_j^{[m]}
\end{equation}

We define the momentum-based selection criterion $\tilde{S}_j^{[m]}$ by subtracting the momentum term $v_j^{[m]}$ scaled with a strength parameter $\beta \ge 0$ and the standard deviation of the current losses $\sigma_S$:

\begin{equation}
    \tilde{S}_j^{[m]} = S_j^{[m]} - \beta \cdot \sigma_S \cdot v_j^{[m]}
\end{equation}

The algorithm then selects the feature minimizing the momentum-based loss $\tilde{S}_j^{[m]}$. Thus, the mechanism incentivizes the selection of features, which consistently reduce the loss over recent iterations, thereby potentially leading to stabilize the selection process.

The integration of the advanced feature selection methods Top-K and Momentum, as well as the flooding regularization into the CWB algorithm is formalized in Algorithm \ref{alg:adv_cwb}. It is possible to derive modularized versions of the algorithm by setting the respective method parameters to ineffectiveness levels: Flood level $b=0$, Momentum strength $\beta=0$, Top-K $K=1$)

\begin{algorithm}
\caption{Advanced CWB with Top-K, Momentum, and Flooding}
\label{alg:adv_cwb}
\begin{algorithmic}[1]
\State \textbf{Initialize:} $\hat{f}^{[0]}(\mathbf{x}) \equiv \bar{y}$, Momentum $\mathbf{v} = \mathbf{0} \in \mathbb{R}^p$.
\State \textbf{Hyperparameters:} Rate $\nu$, Iterations $M_{stop}$, Top-K $K$, Momentum decay $\alpha$, Strength $\beta$, Flood level $b$.
\For{$m = 1$ to $M_{stop}$}
    \State \textbf{Calculate Empirical Risk:} $\mathcal{R}_{emp} = \frac{1}{n}\sum (y_i - \hat{f}^{[m-1]}(\mathbf{x}_i))^2$
    \State \textbf{Compute pseudo-residuals (Flooding):} 
    \If{$\mathcal{R}_{emp} > b$}
        \State $u_i^{[m]} = y_i - \hat{f}^{[m-1]}(\mathbf{x}_i)$ \quad (Gradient Descent)
    \Else
        \State $u_i^{[m]} = -(y_i - \hat{f}^{[m-1]}(\mathbf{x}_i))$ \quad (Gradient Ascent)
    \EndIf
    
    \State \textbf{Fit base learners \& Compute Loss:}
    \For{$j = 1$ to $p$}
        \State Fit $h_j$ to $\mathbf{u}^{[m]}$, compute loss $S_j^{[m]}$.
    \EndFor
    
    \State \textbf{Apply Momentum (if $\beta > 0$):}
    \For{$j = 1$ to $p$}
        \State Compute normalized gain $g_j^{[m]}$.
        \State Update history: $v_j^{[m]} = \alpha v_j^{[m-1]} + g_j^{[m]}$.
        \State Adjust loss: $\tilde{S}_j^{[m]} = S_j^{[m]} - \beta \cdot \sigma_S \cdot v_j^{[m]}$.
    \EndFor
    
    \State \textbf{Select Component (Top-K):}
    \State Identify indices $\mathcal{J}_K$ of the $K$ smallest values in $\tilde{\mathbf{S}}^{[m]}$.
    \State Sample $j^*$ from $\mathcal{J}_K$ using rank-weighted probabilities.
    
    \State \textbf{Update model:}
    \State $\hat{f}^{[m]}(\mathbf{x}) = \hat{f}^{[m-1]}(\mathbf{x}) + \nu \cdot \hat{h}_{j^*}(x^{(j^*)})$
\EndFor
\State \textbf{Return} $\hat{f}^{[M_{stop}]}(\mathbf{x})$
\end{algorithmic}
\end{algorithm}








\newpage
\section{Experimental Design}

\subsection{Data Generation and Preprocessing}
\subsubsection{Synthetic Data Scenarios}
\paragraph{Baseline Signals} 
% The 4 types: Linear, Poly, Step, High-Freq
\paragraph{Deviation Scenarios} 
% Noise magnitude, Dimensionality, Correlation structure
\paragraph{Mixed Signal Compilation} 
% The complex scenario for competing learners
\subsubsection{Real World Datasets}
% Diabetes, Bodyfat, Riboflavin descriptions
\subsubsection{Preprocessing and Scaling}
% Standard scaling, Train/Test splitting logic
\subsection{Study Phases}
\subsubsection{Phase I: Matching Base Learners}
% Linear vs Linear, Spline vs Sine, etc.
\subsubsection{Phase II: Robustness Analysis}
% Testing the deviations (Noise, Correlation)
\subsubsection{Phase III: Competing Base Learners}
% The mixed signal experiment and Real Data setup

\subsection{Evaluation Methodology}
\subsubsection{Hyperparameter Tuning and Grid Search}
\subsubsection{Performance Metrics}
% MSE, Test Loss
\subsubsection{Statistical Significance Testing}
% Paired t-tests methodology

\newpage
\section{Results}

\subsection{Analysis of Synthetic Scenarios}

Consequently, the "floating" mechanism in our CWB implementation does not manifest as a random walk driven by gradient noise, but rather as a result of the deterministic oscillating dynamics between gradient descent and gradient ascent around the flood level $b$. When combined with Top-K selection, this facilitates a "controlled exploration" of the loss landscape, preventing the optimization trajectory from becoming trapped in greedy local optima without compromising the structural stability required for consistent feature selection.

\subsubsection{Performance on Baseline Signals}
% Does CWB fit the correct functions?
\subsubsection{Impact of Flooding and Stochastic Selection}
% Results from the "Deviations" (High noise, etc.)
\subsubsection{Competing Learners on Mixed Signals}
% Does the model correctly identify that it needs Splines for some features and Linear for others?

\subsection{Analysis of Real World Data}
\subsubsection{Benchmark Results on Standard Datasets}
% Bodyfat and Diabetes results
\subsubsection{High-Dimensional Case Study: Riboflavin}
% Specific analysis for the p>>n case

\newpage
\section{Discussion}
\subsection{Interpretation of Generalization Capabilities}
\subsection{The Trade-off between Selection Stability and Predictive Performance}
\subsection{Limitations}

\newpage
\section{Summary and Conclusion}

\newpage
\addcontentsline{toc}{section}{Bibliography}
\renewcommand\refname{Bibliography} 
\bibliographystyle{plainnat}
\bibliography{references}

\newpage
\appendix 
\section{Additional Figures and Tables} 

\end{document}